\documentclass[12pt]{article}
\usepackage[utf8]{inputenc}
\usepackage[T1]{fontenc}
\usepackage{hyperref}
\title{Te'un Language Resources for K-12 Curriculum Development}
\author{Compiled by Claude Sonnet 5 (claude-sonnet-5) for Andr\'as Kornai}
\date{}
\begin{document}
\maketitle

\section{Introduction}
Te'un (ISO 639-3 tve) is an Austronesian language of Maluku, Indonesia. Its speakers originally lived on Teun Island, and they were moved to Seram because of volcanic activity. Sources conflict on its status: it is classed as extinct or nearly so, but one source (Wurm 2007) calls it endangered. Younger generations have shifted to Ambonese Malay and Indonesian. Very little was found, and this report says so plainly.

\section{Primary Educational Resources}
No primers, textbooks or school programmes were found in the search results.

\section{Linguistic and Reference Materials}
\begin{itemize}
\item Wikipedia article on the Teun language: \url{https://en.wikipedia.org/wiki/Teun_language}
\item Endangered Languages catalogue entry: \url{https://www.endangeredlanguages.com/lang/2792?hl=zh-cn}
\end{itemize}

\section{Religious and Community Resources}
\begin{itemize}
\item Joshua Project, Te'un in Indonesia (people group profile): \url{https://joshuaproject.net/people_groups/15452}
\end{itemize}

\section{Digital and Audio Resources}
No audio, video or corpus resources were found in the search results.

\section{Curriculum Development Context}
A heritage curriculum would have to be built almost from scratch. It would draw on whatever archival wordlists exist, which this search did not locate, and on work with the descendant community on Seram. Until such material is found, only cultural and historical heritage content can be documented responsibly.
\end{document}